Это короткий \textbf{учебный образец} того, как зафиксировать ожидание,
заметить ошибку, проверить её в GDB и оформить вывод. Задание не
заменяет лабораторные семинара №4. Рядом опубликован полностью
заполненный пример отчёта: используйте его как ориентир по структуре, а
в своих отчётах приводите собственные команды и наблюдения. Понадобятся
терминал Linux, GCC, GDB и текстовый редактор. В этой демонстрации ввод
не требуется: значения заданы прямо в коде.

\hypertarget{ux43fux43eux440ux44fux434ux43eux43a-ux441ux434ux430ux447ux438-ux434ux435ux43cux43eux43dux441ux442ux440ux430ux446ux438ux43eux43dux43dux43eux439-ux440ux430ux431ux43eux442ux44b}{%
\section{Порядок сдачи демонстрационной
работы}\label{ux43fux43eux440ux44fux434ux43eux43a-ux441ux434ux430ux447ux438-ux434ux435ux43cux43eux43dux441ux442ux440ux430ux446ux438ux43eux43dux43dux43eux439-ux440ux430ux431ux43eux442ux44b}}

\begin{enumerate}
\def\labelenumi{\arabic{enumi}.}
\tightlist
\item
  \textbf{На занятии} выполните часть работы, запишите команды, расчёты
  и наблюдаемый результат и покажите его преподавателю. Это
  промежуточный показ, а не окончательная сдача.
\item
  \textbf{К следующей лабораторной} завершите работу и подготовьте
  отчёт: что сделали, как проверяли, что ожидали и что наблюдали.
  \textbf{Отправьте преподавателю отчёт в формате PDF или HTML.}
  Сохраните исходники и результаты отдельно для демонстрации.
\item
  \textbf{В начале следующей лабораторной} пройдите короткий тест по
  базовой теории предыдущей работы. Тест обязателен, но его результат не
  является допуском к устной защите.
\item
  \textbf{Во время следующей лабораторной} индивидуально защитите
  предыдущую работу: покажите отчёт и работающий результат, объясните
  свои действия и ответьте на вопросы. Преподаватель сообщит, что нужно
  исправить.
\item
  \textbf{Последняя лабораторная} выполняется, показывается и защищается
  на том же занятии. В его начале проходит тест по предпоследней работе;
  отдельного следующего теста по последней работе не будет.
\end{enumerate}

Для этой демонстрации попробуйте оформить свой небольшой отчёт по той же
схеме. Заполненный образец предназначен для изучения оформления, не для
сдачи вместо собственной работы.

\hypertarget{1-ux43fux440ux435ux434ux441ux43aux430ux436ux438ux442ux435-ux440ux435ux437ux443ux43bux44cux442ux430ux442-ux438-ux43dux430ux43fux438ux448ux438ux442ux435-ux43fux440ux43eux433ux440ux430ux43cux43cux443}{%
\section{1. Предскажите результат и напишите
программу}\label{1-ux43fux440ux435ux434ux441ux43aux430ux436ux438ux442ux435-ux440ux435ux437ux443ux43bux44cux442ux430ux442-ux438-ux43dux430ux43fux438ux448ux438ux442ux435-ux43fux440ux43eux433ux440ux430ux43cux43cux443}}

\begin{verbatim}
mkdir -p ~/aiya/demo-sum
cd ~/aiya/demo-sum
nano sum.c
\end{verbatim}

Перепечатайте код в редактор, сохраните через \textbf{Ctrl+O},
\textbf{Enter}, затем выйдите через \textbf{Ctrl+X}:

\begin{verbatim}
#include <stdio.h>

int main(void)
{
    int a = 7;
    int b = 5;
    int result = a - b;
    printf("Сумма: %d\n", result);
    return 0;
}
\end{verbatim}

\textbf{До запуска} запишите в отчёт: задача --- вычислить сумму
\texttt{a\ +\ b}; при \texttt{a\ =\ 7} и \texttt{b\ =\ 5} ожидается
\texttt{12}. Намеренная логическая ошибка в исходнике пока не отмечена:
ваша цель --- обнаружить её по расхождению между ожиданием и
наблюдением.

\hypertarget{2-ux441ux43eux431ux435ux440ux438ux442ux435-ux438-ux441ux440ux430ux432ux43dux438ux442ux435-ux441-ux43eux436ux438ux434ux430ux43dux438ux435ux43c}{%
\section{2. Соберите и сравните с
ожиданием}\label{2-ux441ux43eux431ux435ux440ux438ux442ux435-ux438-ux441ux440ux430ux432ux43dux438ux442ux435-ux441-ux43eux436ux438ux434ux430ux43dux438ux435ux43c}}

\begin{verbatim}
gcc -std=c11 -Wall -Wextra -g -O0 sum.c -o sum
./sum
echo $?
\end{verbatim}

Ожидайте вывод \texttt{Сумма:\ 2}, затем \texttt{0} как код завершения.
Сборка без предупреждений означает, что GCC не обнаружил проблем, о
которых умеет предупредить при этих флагах; она \textbf{не доказывает},
что формула правильная. Нулевой код завершения говорит лишь о том, что
программа сама сообщила об успешном окончании. В отчёте сохраните оба
наблюдения --- \texttt{2} и \texttt{0} --- и сравните первое с
предсказанием \texttt{12}. \texttt{-g} добавляет сведения для отладчика,
\texttt{-O0} делает сопоставление шагов GDB с исходником проще.

\hypertarget{3-ux438ux441ux441ux43bux435ux434ux443ux439ux442ux435-ux440ux430ux441ux445ux43eux436ux434ux435ux43dux438ux435-ux432-gdb}{%
\section{3. Исследуйте расхождение в
GDB}\label{3-ux438ux441ux441ux43bux435ux434ux443ux439ux442ux435-ux440ux430ux441ux445ux43eux436ux434ux435ux43dux438ux435-ux432-gdb}}

В обычном терминале выполните \texttt{gdb\ -q\ ./sum}. После приглашения
\texttt{(gdb)} вводите команды отладчика по очереди:

\begin{verbatim}
break sum.c:7
run
print a
print b
next
print result
quit
\end{verbatim}

\texttt{break\ sum.c:7} остановит программу \textbf{перед} строкой
\texttt{int\ result\ =\ a\ -\ b;} в приведённом коде (если вставляли
строки, номер изменится). \texttt{run} начинает исполнение; после
остановки \texttt{print\ a} покажет \texttt{7}, \texttt{print\ b} ---
\texttt{5}. \texttt{next} выполнит строку вычисления, и только тогда
\texttt{print\ result} покажет \texttt{2}. Числа после знака \texttt{\$}
в ответах GDB --- номера результатов команды \texttt{print}, а не имена
переменных C. \texttt{quit} выходит из отладчика. Запишите в отчёт
наблюдения и объясните, почему получилась двойка.

\hypertarget{4-ux438ux441ux43fux440ux430ux432ux44cux442ux435-ux43aux43eux434-ux438-ux43fux440ux43eux432ux435ux440ux44cux442ux435-ux441ux43dux43eux432ux430}{%
\section{4. Исправьте код и проверьте
снова}\label{4-ux438ux441ux43fux440ux430ux432ux44cux442ux435-ux43aux43eux434-ux438-ux43fux440ux43eux432ux435ux440ux44cux442ux435-ux441ux43dux43eux432ux430}}

Откройте \texttt{nano\ sum.c}, замените \texttt{a\ -\ b} на
\textbf{\texttt{a\ +\ b}}, сохраните файл и снова соберите его той же
командой GCC. Повторный запуск \texttt{./sum} даст \texttt{Сумма:\ 12}.
Запись исходника на диск \textbf{не меняет старый исполняемый файл}: без
пересборки вы по-прежнему увидите прежний результат.

Для двух дополнительных проверок по очереди измените числа в исходнике,
каждый раз пересоберите и запустите программу. Запишите ожидаемый и
наблюдаемый результаты:

\begin{longtable}[]{@{}
  >{\raggedleft\arraybackslash}p{(\columnwidth - 4\tabcolsep) * \real{0.1600}}
  >{\raggedleft\arraybackslash}p{(\columnwidth - 4\tabcolsep) * \real{0.1600}}
  >{\raggedleft\arraybackslash}p{(\columnwidth - 4\tabcolsep) * \real{0.6400}}@{}}
\toprule
\begin{minipage}[b]{\linewidth}\raggedleft
\texttt{a}
\end{minipage} & \begin{minipage}[b]{\linewidth}\raggedleft
\texttt{b}
\end{minipage} & \begin{minipage}[b]{\linewidth}\raggedleft
Ожидаемая сумма
\end{minipage} \\
\midrule
\endhead
\texttt{0} & \texttt{5} & \texttt{5} \\
\texttt{-3} & \texttt{5} & \texttt{2} \\
\bottomrule
\end{longtable}

После проверки верните в файл \texttt{a\ =\ 7} и \texttt{b\ =\ 5} и
пересоберите. Сохраните файл \texttt{sum.c} отдельно от отчёта для
показа. Сравните свой отчёт с заполненным образцом рядом с методичкой: в
нём есть цель, команды, ожидания, наблюдения, найденная ошибка,
повторная проверка и вывод.
